%% =====================================================================
%%  T-21-05  Strike the Shepherd
%%  -------------------------------------------------------------------
%%  One idea per file. Core is mandatory. Slots in canonical order.
%%  Max 700 words. No layout commands. No filler. New source material
%%  for this entry goes in sources/B3/T-21-05.tex -- never by
%%  rewriting this file (docs/RULES.md R0.1).
%%  =====================================================================
%% @kind: topic
%% @id: T-21-05
%% @title: Strike the Shepherd
%% @subtitle: Trouble has an author --- find them, and the flock disperses
%% @chapter: c21
%% @order: 50
%% @status: stable
%% @sources: B3
%% @updated: 2026-09-19

\topic{T-21-05}{Strike the Shepherd}{Trouble has an author --- find them, and the flock disperses}

\begin{core}
Trouble in a group can usually be traced to a single strong individual --- the stirrer, the malcontent whose grievances give everyone else permission. Negotiating with them is wasted: they are untreatable, and delay lets their influence compound. The law is to neutralise the source --- isolate, banish, or remove --- and the sheep will scatter, because the flock was following a person, not a programme.
\end{core}
\begin{mechanism}
B3's judgment is surgical about the diagnosis: beneath every surface of group goodwill there can be one agent keeping dissatisfaction alive, and the group's behaviour collapses quickly without them. The mechanism is narrative supply: the agitator writes the story (who is wronged, who is to blame, what happens next), and the mass supplies only assent --- remove the author and the story has no pen. Athens' ostracism is the institutional form: ten thousand citizens voting to exile the one man whose ambition the city could not otherwise survive. The law's conditions are strict, and stated in the source: strike early, before the trouble multiplies; do not negotiate, because the stirrer's currency is the grievance and negotiation mints more of it; and make the removal about the behaviour's effect, never a public trial of ideas --- a tried agitator becomes an awarded one. The reversal is the guard-rail: strike the wrong target --- or the right target at the wrong moment, turning them into a martyr --- and you have created the author the room lacked.
\end{mechanism}
\begin{conditions}
Strongest in bounded groups with one authority and a identifiable author --- teams, congregations, boards, factions. It weakens where malcontent is genuinely leaderless (the authorless version is culture, and culture cannot be exiled), where the stirrer is load-bearing (talent plus poison is a genuinely hard call), and in public spheres where removal manufactures the martyrdom that sustains the movement.
\end{conditions}
\begin{application}
  \item Diagnose before acting: map who feeds the narrative, who arranges the meetings, who profits from the mood. One name, or you are not done.
  \item Act early. The cheap moment is before the author has an audience; after, every option is expensive.
  \item Remove, don't debate: isolation from the audience, reassignment, exit --- the author must lose the room, not win a forum.
  \item Keep it procedural and quiet. The flock should notice the silence afterwards, not the trial.
  \item Re-map quarterly. Authors are replaced; the replacement is often quieter and better hidden.
\end{application}
\begin{feedback}
  \item Working: the mood deflates faster than the merits explain. The narrative lost its narrator.
  \item Working: no one applies for the vacant authorship. The removal read as maintenance, not persecution.
  \item Failing: a new author surfaces in weeks, better at it. You removed a person; the grievance was the soil.
  \item Failing: the removal makes you the story. Martyrdom achieved --- the flock now has a flag and a founder.
\end{feedback}
\begin{failure}
The tool's abuse is the surveillance state in miniature: a ruler who sees shepherds everywhere manufactures them, and a group that watches its members be removed for murmuring learns to hide every real conversation. It also selects the ambitious into silence rather than agreement, and it optimises for quiet, not health --- a group too well-shepherded stops telling its own leadership what is wrong (T-26-01's mirror entry holds the counterweight).
\end{failure}
\begin{countermeasures}
  \item If your organisation's discontent has a named author, ask the author's question first: what part of the grievance is true? Authority that only removes is authority that has stopped listening.
  \item Refuse the authorship yourself: grievances voiced on the record through open channels are protection, not sedition. The anonymous script is what gets people struck.
  \item Watch leaders who narrate all dissent as one person's malice. The shepherd-strike is the standard tool of the insecure throne (T-26-01).
  \item When you are struck at as the shepherd, go explicit: publish the grievance, the record, the proposal. Martyrs are manufactured from silence.
\end{countermeasures}

\sources{B3}
\seesources
\seealso{T-21-06, T-26-01, T-12-01}
